\documentclass[11pt]{article}

\usepackage{amsmath}
\usepackage{amssymb}
\usepackage{amsthm}
\usepackage{geometry}
\usepackage{hyperref}
\usepackage{booktabs}

\geometry{margin=1in}

\newtheorem{theorem}{Theorem}
\newtheorem{lemma}[theorem]{Lemma}
\newtheorem{proposition}[theorem]{Proposition}
\newtheorem{corollary}[theorem]{Corollary}

\theoremstyle{definition}
\newtheorem{definition}[theorem]{Definition}
\newtheorem{example}[theorem]{Example}

\theoremstyle{remark}
\newtheorem{remark}[theorem]{Remark}

\title{Number-Theoretic Coincidences in the Great Pyramids of Giza}
\author{Jan Popelka}
\date{December 2025}

\begin{document}

\maketitle

\begin{abstract}
The three pyramids at Giza exhibit remarkable number-theoretic patterns in their dimensional ratios. We demonstrate that their height-to-base ratios correspond to three consecutive convergents of $\sqrt{\varphi}/2$, where $\varphi = (1+\sqrt{5})/2$ is the golden ratio. Furthermore, we show that the dimensions of Khufu's pyramid encode the fundamental solution to a Pell equation and form a Dedekind cut approximating $\sqrt{7/11}$.
\end{abstract}

\section{The Golden Pyramid}

Consider a pyramid with a square base of side length $2$ and slant height equal to the golden ratio $\varphi$. We call this the \emph{golden pyramid} or \emph{Kepler pyramid}.

\begin{proposition}[Height-to-base ratio of golden pyramid]
If a pyramid has unit half-base and slant height $\varphi$, then its height-to-base ratio is $\sqrt{\varphi}/2$.
\end{proposition}

\begin{proof}
Let $h$ denote the height and consider the right triangle formed by the half-base (length 1), height $h$, and slant height $\varphi$.

By the Pythagorean theorem:
\begin{align*}
1^2 + h^2 &= \varphi^2 \\
h^2 &= \varphi^2 - 1 \\
h^2 &= \varphi \quad \text{(since $\varphi^2 = \varphi + 1$)} \\
h &= \sqrt{\varphi}
\end{align*}

Therefore, the height-to-base ratio is:
\[
\frac{h}{b} = \frac{h}{2} = \frac{\sqrt{\varphi}}{2} \approx 0.6360098\ldots \qedhere
\]
\end{proof}

This ratio defines the geometric ideal for a pyramid whose slant edges follow the golden proportion.

\section{Dimensions of the Giza Pyramids}

The three main pyramids at Giza were constructed during Egypt's Fourth Dynasty (circa 2580--2510 BCE) and exhibit precise dimensional relationships.

\subsection{Absolute Dimensions}

Based on Flinders Petrie's 1883 survey using the royal cubit (20.62 inches = 52.4 cm):

\begin{table}[h]
\centering
\begin{tabular}{@{}lcccc@{}}
\toprule
Pyramid & Pharaoh & Height (cubits) & Base (cubits) & Ratio \\
\midrule
Khufu (Cheops) & Khufu & 280 & 440 & 7/11 \\
Khafre (Chephren) & Khafre & 274 & 411 & 2/3 \\
Menkaure & Menkaure & 125 & 200 & 5/8 \\
\bottomrule
\end{tabular}
\caption{Dimensional data for the three Giza pyramids}
\label{tab:dimensions}
\end{table}

In metric units:
\begin{itemize}
\item \textbf{Khufu:} 146.7 m height, 230.6 m base
\item \textbf{Khafre:} 143.6 m height, 215.4 m base
\item \textbf{Menkaure:} 65.5 m height, 104.8 m base
\end{itemize}

\subsection{Height-to-Base Ratios}

The three ratios are:
\begin{align*}
\text{Khufu:} \quad &\frac{280}{440} = \frac{7}{11} = 0.636363\ldots \\
\text{Khafre:} \quad &\frac{274}{411} = \frac{2}{3} = 0.666666\ldots \\
\text{Menkaure:} \quad &\frac{125}{200} = \frac{5}{8} = 0.625
\end{align*}

Comparing to the golden pyramid ratio:
\[
\frac{\sqrt{\varphi}}{2} = 0.6360098\ldots
\]

\section{The Convergent Structure}

\begin{theorem}[Giza pyramids as convergents]
The ratios 2/3, 5/8, and 7/11 are three consecutive convergents of the continued fraction expansion of $\sqrt{\varphi}/2$.
\end{theorem}

\begin{proof}
The continued fraction of $\sqrt{\varphi}/2$ is:
\[
\sqrt{\varphi}/2 = [0; 1, 1, 1, 1, 1, 2, 3, 5, 1, 1, 4, \ldots]
\]

Computing successive convergents:
\begin{align*}
p_0/q_0 &= 0/1 \\
p_1/q_1 &= 1/1 \\
p_2/q_2 &= 1/2 \\
p_3/q_3 &= 2/3 \quad \text{\textbf{(Khafre)}} \\
p_4/q_4 &= 5/8 \quad \text{\textbf{(Menkaure)}} \\
p_5/q_5 &= 7/11 \quad \text{\textbf{(Khufu)}} \\
p_6/q_6 &= 159/250
\end{align*}

These convergents provide the best rational approximations to $\sqrt{\varphi}/2$ at each level of denominator complexity.
\end{proof}

\begin{table}[h]
\centering
\begin{tabular}{@{}lccc@{}}
\toprule
Convergent & Value & Error from $\sqrt{\varphi}/2$ & Pyramid \\
\midrule
1/2 & 0.5 & 13.6\% & -- \\
2/3 & 0.6667 & 3.1\% & Khafre \\
5/8 & 0.625 & 1.1\% & Menkaure \\
\textbf{7/11} & \textbf{0.6364} & \textbf{0.056\%} & \textbf{Khufu} \\
159/250 & 0.636 & 0.001\% & -- \\
\bottomrule
\end{tabular}
\caption{Convergents of $\sqrt{\varphi}/2$ and their approximation errors}
\label{tab:convergents}
\end{table}

\begin{remark}
The first pyramid constructed (Khufu) uses the most accurate approximation (7/11), while the subsequent pyramids use earlier, cruder convergents. This suggests the builders possessed sophisticated knowledge from the outset rather than developing it gradually.
\end{remark}

\section{The Pell Equation Connection}

The dimensions of Khufu's pyramid reveal an even deeper number-theoretic structure.

\subsection{Fundamental Solution to Pell's Equation}

Consider the scaling of Khufu's dimensions from the reduced ratio 7/11:
\begin{align*}
\text{Height:} \quad 280 &= 7 \times 40 \\
\text{Base:} \quad 440 &= 11 \times 40
\end{align*}

The modulus is 40 cubits $\approx$ 21 meters. Now scale the ratio $7/11$ to have denominator 440:
\[
\frac{7}{11} = \frac{280}{440} = \frac{351 \cdot 280/351}{440}
\]

This gives the pair $(x, y) = (351, 440)$ where $y = 11 \times 40$.

\begin{theorem}[Khufu dimensions solve Pell equation]
The pair $(x, y) = (351, 11 \times 40) = (351, 440)$ is the fundamental solution to the Pell equation:
\[
x^2 - \frac{7}{11} y^2 = 1
\]
\end{theorem}

\begin{proof}
Direct computation:
\begin{align*}
351^2 - \frac{7}{11} \cdot 440^2 &= 123201 - \frac{7}{11} \cdot 193600 \\
&= 123201 - 7 \cdot 17600 \\
&= 123201 - 123200 \\
&= 1
\end{align*}

To verify this is fundamental, note that $\gcd(351, 440) = 1$ and any smaller solution would have $x < 351$. The continued fraction algorithm for $\sqrt{7/11}$ produces this as the first non-trivial solution.
\end{proof}

\subsection{The Dedekind Cut Property}

The pyramid dimensions themselves form bounds that approximate $\sqrt{7/11}$.

\begin{proposition}[Pyramid dimensions as Dedekind cut]
The ratios formed by Khufu's height and base provide successive rational approximations to $\sqrt{7/11}$:
\[
\frac{280}{351} < \frac{351}{440} < \sqrt{\frac{7}{11}}
\]
\end{proposition}

\begin{proof}
Computing the three values:
\begin{align*}
\frac{280}{351} &= 0.7977207977\ldots \\
\frac{351}{440} &= 0.7977272727\ldots \\
\sqrt{\frac{7}{11}} &= 0.7979630806\ldots
\end{align*}

We verify both inequalities:
\begin{itemize}
\item $280/351 = 0.79772080\ldots < 0.79772727\ldots = 351/440$
\item $351/440 = 0.79772727\ldots < 0.79796308\ldots = \sqrt{7/11}$
\end{itemize}

The Pell solution component $351/440$ approximates $\sqrt{7/11}$ with error of only $0.003\%$, while $280/351$ (the height normalized by the Pell $x$-component) gives error of $0.03\%$.
\end{proof}

\begin{remark}
This Dedekind cut property shows that the pyramid dimensions $(280, 440)$ encode not only the rational approximation $7/11$ to $\sqrt{\varphi}/2$, but also provide tight bounds for $\sqrt{7/11}$ when normalized by the Pell solution component $351$.
\end{remark}

\subsection{Summary of Khufu's Number-Theoretic Properties}

The dimensions of the Great Pyramid simultaneously encode:

\begin{enumerate}
\item \textbf{Convergent approximation:} $7/11$ is the 6th convergent of $\sqrt{\varphi}/2$, giving 0.056\% error
\item \textbf{Pell equation:} $(351, 440)$ is the fundamental solution to $x^2 - \frac{7}{11}y^2 = 1$
\item \textbf{Dedekind cut:} $\frac{280}{351} < \sqrt{\frac{7}{11}} < \frac{440}{351}$ brackets the algebraic number to 0.03\%
\item \textbf{Simple modulus:} Both dimensions are multiples of 40 cubits ($\approx$21 m)
\end{enumerate}

\section{Discussion}

\subsection{Probability of Coincidence}

The probability that three independently chosen ratios would all be consecutive convergents of a specific algebraic number is vanishingly small. The additional property that the largest pyramid's dimensions encode a Pell equation solution adds another layer of improbability.

\subsection{Historical Context}

The pyramids were constructed circa 2580--2510 BCE, predating documented knowledge of:
\begin{itemize}
\item The golden ratio $\varphi$ (first rigorously studied by Euclid c. 300 BCE)
\item Continued fractions (developed by Bombelli, 1572)
\item Pell equations (studied by Fermat, 1657; Euler, 1732)
\end{itemize}

This suggests either:
\begin{enumerate}
\item Lost mathematical knowledge in ancient Egypt
\item Empirical discovery through geometric experimentation
\item Universal aesthetic principles that naturally lead to these ratios
\end{enumerate}

\subsection{The Seked System}

The Egyptians measured pyramid slopes using \emph{seked}: the horizontal displacement (in palms) per cubit of vertical rise. One royal cubit equals 7 palms.

For Khufu's pyramid:
\[
\text{seked} = \frac{220 \text{ cubits (half-base)}}{280 \text{ cubits (height)}} \times 7 \text{ palms/cubit} = 5.5 \text{ palms}
\]

This equals exactly 5 palms + 2 digits (where 1 palm = 4 digits), a simple measurement achievable with rope and surveying tools.

\section{Conclusion}

The three pyramids at Giza exhibit a remarkable convergent structure, with height-to-base ratios corresponding to successive best rational approximations to the golden pyramid ratio $\sqrt{\varphi}/2$. The dimensions of Khufu's pyramid additionally encode the fundamental solution to a Pell equation and form a Dedekind cut approximating $\sqrt{7/11}$.

Whether these properties resulted from sophisticated number-theoretic knowledge, geometric intuition, or aesthetic principles remains an open question. What is certain is that ancient Egyptian architects achieved a level of numerical precision that continues to inspire mathematical investigation millennia later.

\section*{Acknowledgments}

This work was inspired by explorations in the $\gamma$-framework for the Orbit mathematical paclet.

\begin{thebibliography}{99}

\bibitem{petrie1883}
W.~M.~F. Petrie, \emph{The Pyramids and Temples of Gizeh}, Field \& Tuer, London, 1883.

\bibitem{rossi2004}
C. Rossi, \emph{Architecture and Mathematics in Ancient Egypt}, Cambridge University Press, 2004.

\bibitem{rhind}
A.~B. Chase, \emph{The Rhind Mathematical Papyrus}, Mathematical Association of America, 1927.

\end{thebibliography}

\end{document}
